% Gauss, Werke X/1 (1917), native58–59, printed51–52.
\SourceLeaf{058}\NoSourceHead
\SourceGroup{VII}
\HistoricTitle{BEWEIS EINES SATZES AUS DER HÖHERN ARITHMETIK ZUR THEORIE DER BIQUADRATISCHEN RESTE GEHÖRIG.}
\HistoricCitation{[Aus Handbuch 21, Bg, S. 5, 6.]}
\par\indent Es sei $a$ eine ganze positive oder negative Zahl von der Form $4m+1$, $n$ eine ganze positive Zahl, die mit $a$ keinen Theiler gemein habe. $f$ bedeute den Complexus der Zahlen $1,2,3,\ldots,n$, $f'$ die ungeraden darunter, $f''$ die geraden; ferner bedeute $g$ die Zahlen $\tfrac12f''$, oder alle Zahlen $1,2,3,\ldots,\tfrac12n$ oder $1,2,3,\ldots,\tfrac12(n-1)$, jenachdem $n$ gerade oder ungerade, $h$ die übrigen $\tfrac12n+1$, $\tfrac12n+2,\ldots,n$ oder $\tfrac12(n+1)$, $\tfrac12(n+3),\ldots,n$. Man hat sodann
\[
\tfrac14n(a-1)=\sum\left[\frac{fa}{4n}+\frac14\right]+\sum\left[\frac{fa}{4n}+\frac12\right]
\]
\[
\begin{aligned}
\tfrac14n(a-1)&=\sum\left[\frac{ga}{2n}+\frac12\right]+\sum\left[\frac{ha}{2n}\right]
=\sum\left[\frac{ga}{2n}+\frac12\right]-\sum\left[\frac{ga}{2n}\right]+\sum\left[\frac{fa}{2n}\right]\\
&=\sum\left[\frac{ga}{2n}+\frac12\right]-\sum\left[\frac{ga}{2n}\right]+\sum\left[\frac{fa}{4n}\right]+\sum\left[\frac{fa}{4n}+\frac12\right]\\
&=\sum\left[\frac{ga}{2n}+\frac12\right]+\sum\left[\frac{f'a}{4n}\right]+\sum\left[\frac{fa}{4n}+\frac12\right],
\end{aligned}
\]
\noindent folglich
\[
\sum\left[\frac{fa}{4n}+\frac14\right]=\sum\left[\frac{ga}{2n}+\frac12\right]+\sum\left[\frac{f'a}{4n}\right]
\]
\noindent und
\[
\sum\left[\frac{fa}{4n}+\frac12\right]-\sum\left[\frac{fa}{4n}+\frac14\right]=\sum\left[\frac{f'a}{4n}+\frac12\right]-\sum\left[\frac{f'a}{4n}\right].
\]
\noindent Nun aber ist
\[
\left[\frac{ha}{4n}+\frac12\right]-\left[\frac{ha}{4n}+\frac14\right]=\left[\frac{(n-h)a}{4n}+\frac12\right]-\left[\frac{(n-h)a}{4n}+\frac14\right],
\]
\par\vfill\noindent\hfill7*\par
\clearpage

\SourceLeaf{059}\ArithmeticHead{52}{L}{ARITHMETIK.\quad NACHLASS.}
\noindent also
\[
\sum\left[\frac{fa}{4n}+\frac12\right]-\sum\left[\frac{fa}{4n}+\frac14\right]=2\sum\left[\frac{ga}{4n}+\frac12\right]-2\sum\left[\frac{ga}{4n}+\frac14\right],
\]
\noindent wenn $n$ ungerade,
\[
=2\sum\left[\frac{ga}{4n}+\frac12\right]-2\sum\left[\frac{ga}{4n}+\frac14\right]-\left[\frac a8+\frac12\right]+\left[\frac a8+\frac14\right],
\]
\noindent wenn $n$ gerade.

Das Theorem kann also auch so ausgedrückt werden:\newline
»Die Summe aller
\[
\left[\frac{f'a}{4n}+\frac12\right]-\left[\frac{f'a}{4n}\right],
\]
\noindent wo $f'$ alle ungeraden Zahlen bis $n$ inclusive bedeutet, ist gleich der doppelten Summe aller
\[
\left[\frac{ga}{4n}+\frac12\right]-\left[\frac{ga}{4n}+\frac14\right],
\]
\noindent wenn $g$ alle Zahlen bis $\tfrac12n$ bedeutet, nur dass für $g=\tfrac12n$ nur der halbe Werth gezählt werden muss«.
\ClosingRule
\begin{center}\textbf{BEMERKUNG.}\end{center}
\begingroup\small
\par\indent\textsc{Gauss} hat die Gelegenheiten, wo der hier mitgeteilte Satz in der Theorie der biquadratischen Reste verwendbar ist, nicht näher angedeutet; vermutlich hat er ihm bei Umformungen gedient ähnlich denjenigen, wie sie in der \textit{Commentatio secunda}, art. 73—75, Werke II, S. 143 ff. und verwandten Aufzeichnungen auftreten.
\begin{flushright}\textsc{Bachmann.}\end{flushright}
\ClosingRule\clearpage\endgroup
